\section{Data, target, and evaluation
design}\label{data-target-and-evaluation-design}

\begin{figure}
\centering
\includegraphics{figures/fig1_research_design.pdf}
\caption{Research design, chronological split, and delayed information
flow.}\label{fig:figure-1}
\end{figure}

\subsection{3.1 NBS state-commodity development
panel}\label{nbs-state-commodity-development-panel}

The development panel contains monthly retail food prices from March
2017 to April 2024. It covers 37 state units, including the Federal
Capital Territory, and 43 food items. After harmonisation, the balanced
key structure contains 136,826 state-item-month observations. NBS
reports state averages based on field collection across the federation
and applies internal verification procedures. The modelling panel also
includes month-t food-price dynamics and selected macroeconomic
indicators, including premium motor spirit prices, exchange rates, and
transport fares.

Several isolated price records display implausible one-month unit or
decimal spikes. The primary upper-tail target remains quantile-based,
which is less sensitive than squared-error objectives, and a
training-based 0.5/99.5 percentile winsorisation is reported as a
robustness check. No outlier rule is estimated using the test period.

\subsection{3.2 Same-source NBS national
continuation}\label{same-source-nbs-national-continuation}

For temporal persistence, later NBS releases were harmonised to 28
national item series through April 2026. Product-unit matches were
retained only when the old and new labels were sufficiently comparable.
Exact calendar gaps were preserved, and a three-month target was formed
only when the exact t+3 observation existed. The post-2022 evaluation
contains 638 item-month forecasts across 23 forecast-origin months.

\subsection{3.3 External WFP market
panel}\label{external-wfp-market-panel}

The external panel uses observational WFP retail prices in Nigerian
naira. Eligible market-commodity-unit series require at least 18
pre-2022 targets, at least two validation targets, and at least six
post-2022 targets. The resulting panel contains 366 series, 15 markets,
25 commodities, 7,123 post-2022 forecasts, and 38 forecast-origin
months. Post-2022 coverage is concentrated in Borno and Yobe, so this is
a cross-source and market-level transport test, not a nationally
representative replication.

\textbf{Table 1. Empirical datasets and roles}

\begin{longtable}[]{@{}
  >{\raggedright\arraybackslash}p{(\columnwidth - 8\tabcolsep) * \real{0.1630}}
  >{\raggedright\arraybackslash}p{(\columnwidth - 8\tabcolsep) * \real{0.1630}}
  >{\raggedright\arraybackslash}p{(\columnwidth - 8\tabcolsep) * \real{0.1630}}
  >{\raggedright\arraybackslash}p{(\columnwidth - 8\tabcolsep) * \real{0.3261}}
  >{\raggedright\arraybackslash}p{(\columnwidth - 8\tabcolsep) * \real{0.1630}}@{}}
\toprule\noalign{}
\begin{minipage}[b]{\linewidth}\raggedright
\textbf{Dataset}
\end{minipage} & \begin{minipage}[b]{\linewidth}\raggedright
\textbf{Period}
\end{minipage} & \begin{minipage}[b]{\linewidth}\raggedright
\textbf{Units}
\end{minipage} & \begin{minipage}[b]{\linewidth}\raggedright
\textbf{Scale}
\end{minipage} & \begin{minipage}[b]{\linewidth}\raggedright
\textbf{Role}
\end{minipage} \\
\midrule\noalign{}
\endhead
\bottomrule\noalign{}
\endlastfoot
NBS development panel & Mar 2017-Apr 2024 & 37 states/FCT × 43 items;
136,826 rows & State-item-month & Model development and locked 2023
break test \\
NBS national continuation & Through Apr 2026 & 28 harmonised item
series; 638 post-2022 forecasts & National item-month & Same-source
temporal extension and recovery test \\
WFP market observations & Through 2026 & 366 series; 15 markets; 25
commodities; 7,123 post-2022 forecasts & Market-commodity-unit-month &
External source and market-level transport test \\
\end{longtable}

Notes: The WFP external sample is concentrated in Borno and Yobe after
the predeclared eligibility screen. The World Bank Real-Time Food Prices
catalogue is used as a reproducibility reference but its imputed prices
are not treated as independent observational ground truth.

\subsection{3.4 Forecast target}\label{forecast-target}

For state or market unit s, commodity i, and forecast-origin month t,
the three-month forward log price change is

\[
y_{ist}^{(3)}=100\left[\log(P_{is,t+3})-\log(P_{ist})\right].
\]

The principal target is the conditional 90th percentile
\(q_{0.90}(y_{ist}^{(3)}\mid\mathcal{I}_t)\), denoted
\(\mathrm{FIaR}_{0.90}^{(3)}\). Forecasts are issued after month-t
prices and contemporaneous predictors are observed. The outcome attached
to forecast origin t becomes available only after month t+3. Therefore,
at origin t, adaptive updates may use only errors attached to origins
t-3 or earlier.

\subsection{3.5 Chronological splits and inferential
status}\label{chronological-splits-and-inferential-status}

The base quantile LightGBM is trained through December 2020.
Regime-conditioned priors use 2021. Hyperparameters and controller
penalties use 2022. The original development test is January 2023 to
January 2024; continuation and external evaluations extend beyond that
period. The split is chronological throughout, and no random
cross-validation is used.

The Phase 10 controller is nevertheless an exploratory second-stage
backtest. Its architecture was proposed after earlier adaptive models
had already been evaluated on the post-2022 period. The 2023--2026
estimates therefore demonstrate internal consistency and practical
promise, not pristine confirmatory superiority. Bootstrap intervals are
descriptive uncertainty assessments.
